% Oppgaven kan skrives på norsk eller engelsk. Hvis den skrives på engelsk må i så fall overskriftene oppdateres.
% bruk oppsett her
% Merk at 25% av karakteren på oppgaven er struktur og form og vi anbefaler derfor alle om følge denne malen for å unngå unødvendig trekk i karakter.

% Husk å kjøre pdflatex 2 ganger etter bibtex:
% pdflatex filnavn
% bibtex filnavn
% pdflatex filnavn
% pdflatex filnavn

% MERK: det er nok flere pakker som må innstalleres her!
\documentclass[11pt]{article} 
\usepackage[utf8]{inputenc} 
\usepackage{geometry} 
\geometry{a4paper} 
\usepackage{graphicx} 
\usepackage{fancyhdr}
\usepackage[english]{babel} % 'norsk' må endres til 'english' hvis du skal skrive på Engelsk
% Inkluder eventuelle pakker her

% START:  DISSE MÅ ENDRES
\newcommand{\tittel}{TITTEL} % Tittel velger du 
\newcommand{\temanr}{26} % referer til det som er oppgitt på wiki-siden
\newcommand{\tematittel}{Tracking and mass surveillance in the era of Big Data} % referer til det som er oppgitt på wiki-siden 
\newcommand{\kandidatnr}{julihau, olejfri} % komma-separert liste med kandidatnummer for alle forfatter VIKTIG: bruk kandidatnummer og ikke studentnummer siden oppgaven skal være anonym.
\newcommand{\kandidatnrforside}{julihau \\ olejfri}  %  liste med kandidatnummer for alle forfatter - her skal hver forfatter separeres med \\
\newcommand{\antallord}{ORD} % antall ord må fylles inn (trenger ikke være 100% korrekt men skal gi en indikasjon; flere måter texcount, kompiler til PDF og tell, kopier i word og bruk wordcount, +++,  
% Merk: lengden skal være mellom 5000 – 10000 ord 
% Man må påregne trekk i karakter dersom oppgaven er nevneverdig kortere eller lengre enn dette, men antall ord innenfor gitt intervall vil ikke ha betydning på karakter. Her bedømmes det utfra kvalitet og ikke kvantitet. 
% SLUTT:  DISSE MÅ ENDRES


% START: DETTE SKAL IKKE ENDRES
\pagestyle{fancy}
\fancyhf{}
\fancyhead[RO,RE]{\kandidatnr}
\fancyhead[LO,LE]{IN2120 -- midtsemestereksamen 2022}
\fancyfoot[C]{\thepage}
\begin{document} \thispagestyle{empty}
\begin{center}
\normalsize{\textit{Universitetet i Oslo -- IN2120 Informasjonssikkerhet-- 2022 -- midtsemestereksamen}} \\
\vspace{1cm}
\includegraphics[width=0.15\textwidth]{uio_logo.png} \\
\vspace{1cm}
\huge{\tittel} \\
\Large{\temanr} \\
\Large{\tematittel} \\
\vspace{1cm}
\Large{Kandidatnummer:} \\
\large{\kandidatnrforside} \\
\vspace{1cm}
\Large{Antall ord: \antallord} \\
\end{center}
\tableofcontents
\newpage
% SLUTT: DETTE SKAL IKKE ENDRES


% HER KOMMER SELVE TEKSTEN I OPPGAVEN

\section{Abstract}
How have legislations, attitudes towards and the consequences of tracking and mass surveillance changed over the course of time, taking into consideration the recent global pandemic and how data-tracking was used as a countermeasure. 
\\ \\
The practice of mass surveillance was first put to use as a result of the happenings of 9/11 2001 (Sinha, 2013) with the reasoning that it could be used to help prevent future terrorist attacks and the like. It was presented as a necessary evil in order to assist authorities in ensuring the safety of the people, and as such, a new normal emerged. By using massive sets of data collected from its netizens, new systems were created to predict anything from traffic jams to premediated crime (Zuboff, 2015).
\\ \\
The people’s willingness to volunteer their data likely stems from the theory that in order to feel a sense of safety from these frightening prospects, a person would rather sacrifice their own privacy than risk becoming a victim themselves (Zuboff, 2019). This would however turn out to be the beginning of something much bigger, namely the understanding of how much you could actually discern about a person by analyzing their data, and how big companies could profit from doing exactly that (Vargo \& Hopp, 2020).
\\ \\
There have been a few notable cases of tracking, manipulation and surveillance is years following the start of this new paradigm. One example being the Cambridge Analytica scandal, in which it was brought to light that Facebook had handed over personally identifiable information of more than 87 million users to to Cambridge Analytica. It would later be discovered that this information was used to influence the American presidential election of 2016 (Rathi, 2019). What was made clear following scandals such as this was that this new era carried with it the risk of private companies being able to control the masses (Karas, 2002), even to the point of national leaders being elected by their will. 
\\ \\
During the Covid-19 pandemic, multiple forms of mobile applications with the express purpose of tracking the spread of the virus were introduced. Both government-approved applications and ones developed by notorious private companies such as Apple and Google (Kaashoek \& Santillana, 2020) were put to use on an international level (Munzert, et al., 2021). This begs the question; “Where does user privacy begin and end?” How can we be sure that the data we are giving away, with the intent of benefitting the public good, is used in a manner that is effective or at all useful?” Could this be the start of an era where mass data-collection and surveillance are an everyday phenomenon? What data will be standard procedure to collect from users of mobile applications, and how could this information be used in the future? It would seem that it all boils down to the question; “In what ways are we allowing both authorities and private companies to influence our lives by sharing our data?”. 
\\ \\

%Innledningen skal gi en kort introduksjon til tema og motivasjon til oppgaven, og beskrive hva oppgaven skal besvare. Denne vil typisk være rundt en side (men dette er ikke et krav vil variere fra tema til tema).
%Sitering/referanser: 25% av karakteren er god bruk av kilder og referanser. Vi anbefaler bruk av passende sitering i teksten (se under) og vil anbefale at dere bruker bibtex.
% Dere kan bruke hvilken stil som dere ønsker og her har vi brukt en enkel tall-referanse (Word-malen bruker Harvard-stilen) - det er helt opp til dere hvilken stil som brukes.

%Eksempel referanse: \cite{josang21} (denne setningen er for illustrasjon og må fjernes)%

\section{Hovedkapittel}

%Besvarelsen må ha ett eller flere hovedkapitler.  Hvert hovedkapittel kan ha så mange delkapitler som ønskelig (så fremt det gir mening strukturmessig) og skal være nummert som gjort i denne malenn. Fjern/legg til etter behov. Anbefaler klipp-og-lim eksisterende kapittel/delkapittel for å få riktig nummerte overskrifter.

\subsection{Targeted Advertisement}
Due to the increased popularity of personal profiles on social media sites and the like, we find ourselves in a period of mass personalization (Mulvenna, et al., 2000). Many websites and applications like Facebook, Twitter and Tumblr can potentially create a comprehensive list of personal data such as demographic, behavioral, psychographic and lifestyle-data simply by using its intended functions (Russel, 2013) (Arya, et al., 2019). In addition to this mined data, many applications such as WhatsApp, Messenger and TikTok include a statement in their policies, informing us that by using their applications, the user is consenting to their communication being put through textual analysis and searched for keywords. These keywords are then used to personalize what advertisements are shown to the user and other users who display similar textual patterns (Vargo \& Hopp, 2020). 
\\ \\
According to (Trang, 2017), advertising strategies have been changed drastically as a result of the development of social media. Due to their 1.39 million active monthly users, Facebook has become one of the most common online marketing tools with 92\% of social marketing companies having chosen to use the platform due to their comparatively low prices and massive user base.  
\\ \\
In addition to this traditional form of advertisement, Facebook teamed up with IBM on May 6th, 2015 with the purpose of creating better advertising on what was then, and continues to be, the world’s largest social network. By combining Facebook’s targeting technology and IBM’s services for marketers, they created a new form of advertisement system where users’ interactions on collaborating marketers’ websites resulted in advertisements of viewed products or services being generated and put on their Facebook feed. An example of this would be if a user visited company X’s website and looked at a certain product. This viewing data would then be sent by the company to Facebook, and an ad of that product featuring an affiliate link to company X’s website would then later appear on the user’s Facebook feed (Trang, 2017).  
\\ \\
While this practice has been shown to have very high success rates compared to the more traditional methods of advertising, it has brought with it the consequence that selling these massive amounts of data to third parties has become a big business in the 21st century (Palos-Sanchez, et al., 2019). 
\\ \\

\subsection{Manipulation and Recruitment}
We have already discussed why using or selling user data for advertising purposes may be one of the main reasons for websites and applications to collect it, but taking into account the questions we posed earlier, there is another aspect to consider. Namely user data being used by potentially malicious parties wishing to manipulate or recruit the user.
\\ \\ 
As previously mentioned, one notorious example of users can be manipulated as a result of their data being shared is the Cambridge Analytica scandal. Cambridge Analytica was officially an advertising company with both a commercial and a political wing. The political wing was claimed to “combine the predictive data analytics, behavioral sciences, and innovative ad tech into one award winning approach”. The method in practice involved making unique advertisements for many different categories of users which were grouped based on their analyzed traits. These analyzed traits were the results of a regression model that in the words of Professor David Sumpter: “Takes the data we already know about a person, and uses them to predict something we don’t know about him or her” (Sumpter, 2018). The idea behind it all was that different types of people will react more strongly to different forms of advertisement, be it appealing to the viewers sentimentality or by inciting feelings of anger with the opposition. By exposing Facebook users to these tailored ads, they intended to yield better results compared to the more traditional blanket advertisement where the same ad is shown to everyone.  
\\ \\
The effectiveness of this approach has been both exaggerated and underplayed by both the media, CEO Alexander Nix, and the whistleblower and ex-employee Christopher Wylie, but the fact remains that Cambridge Analytica, to an unknown extent, helped republican politician Ted Cruz get the second highest number of votes in the republican primary election of 2016 (Rathi, 2019). Despite the arguments of both sides regarding how effective or ineffective these kinds of regression models are, users of platforms such as Facebook are in fact having their data used in what has been called both “an ethical grey area” and “an attempt to manipulate voters by latching onto their vulnerabilities” (McCausland \& Schnecter, 2018). 
\\ \\
In addition to this form of manipulation, we can also refer to the research done by (Alava, et al., 2017) which speaks of a relatively recent, but growing, phenomenon often referred to as “incitement to radicalization towards violent extremism” or more simply put “violent radicalization”. The situation we are faced with according to their research is that this boom of social media has created new and easier ways for radical groups to identify and approach users that they believe could become allies to their cause. Taking into account that these groups, given sufficient funds could purchase potentially massive amounts of user information to aid them in their recruitment, puts to question whether the premise that the gathering and availability of personal data is in the best interest of public safety. 
\\ \\
Taking what we have already discussed about the subject into account, we can say that by using applications and services like Facebook, Twitter and Snapchat, the user is either knowingly or unknowingly offering their data to the company and whomever they choose to share that data with. This data may be used against the user in the form of targeted advertising, manipulation or potentially tracking and surveillance depending on what data the application chooses to collect. One could argue that a common denominator that allowed, and still allows, these happenings to take place is the culture of mass data collection which, as previously mentioned, has created a thriving market for the commerce and utilization of user data.  
\\
\subsection{Governmental data collection, tracking and surveillance}
Paraphrasing the words of (Cate, 2008), The United States government has long sought data about individuals for a wide variety of important public purposes, but due to the then time and cost-heavy processes yielding mostly useless data, this collecting of data was protected under The Supreme Courts concept of “practical obscurity”. 
\\ \\
The concept stems from the case “U.S. Department of Justice v. Reporters Committee for Freedom of Press” of 1989, and in practice means that while some information is technically accessible to the public, the time and effort needed to acquire this information makes it somewhat protected from mass access. The example of this used in the Supreme court case was that whilst the rap sheets and court records of a person were technically available to any person with the time to dig through a county courthouse’s paper records in search of them, the high cost and low chance of acquiring them meant that the information had little chance of being compiled by the public (Reporters Committee for Freedom of the Press, 2009). 
\\ \\
While this concept was deemed adequate protection of privacy back in 1989, with new technology, the situation has changed. Despite court records and arrest reports being available only as physical copies in most U.S. counties, corporate data brokers are free to collect this information on individuals from all over the country which they can then sell to government agencies, insurance companies and other third parties belonging to the private sector. A notable example of such a company is ChoicePoint which offers aggregate reports of a person’s line of credit collected from multiple public and private databases (Personal Finance, 2019). These reports can in some cases determine whether a person is considered qualified for a job, professional license, or insurance (Testa-Avila, 2008).
\\ \\
While the government does not acquire this data directly, the Supreme Court has refused to restrict the government’s access to data held by third parties which they could do by extending the Fourth Amendment that was originally meant to ensure the right to privacy (Cate, 2008). As it currently stands, companies have common authority over their business records and can therefore consent to a government search or acquisition of their databases even if their users oppose it (Kerr, 2021).
This enforces that by using applications, websites or services, the user allows both private sector parties and authorities access to their personal data which could be used for both surveillance and tracking.
\\
\subsection{The COVID-19 pandemic}
The question of whether or not the handling and management of user data is a violation of user privacy has been the subject of many studies in this era of Big Data (Buchanan, et al., 2006). Something that brought more attention to the issue was however the introduction of the COVID-19 related tracking of citizens.
\\ \\
Since the beginning of the virus-outbreak in 2020, the pandemic caused many systematic changes in organizational structures and society as a whole. Surveillance and control in the form of collecting data about location and interactions for its citizens was advised worldwide, which caused fear for the thought of a future where pandemic countermeasures could systematically invade user privacy (Ribeiro-Navarrete, et al., 2021). 
\\ \\
After the recognition of the second wave of the COVID-19-pandemic, increasingly more countries chose to implement large-scale technology-based tracking measures in order to monitor and prevent the spread of the virus across the population. Despite the governments’ efforts, these technology-based tracking applications ultimately experienced a global failure to convince the public due to these applications’ perceived risk of infringing individual liberties and were met with very low levels of trust from the population (Georgieva, et al., 2021). 
\\ \\
The findings in the article by (Hargittau, et al., 2020) describe the public distrust towards contact-tracing apps as being the product of a decade of media-covered security breaches and other violations of users’ privacy expectations with one notable example being ”Edward Snowden’s revelations of mass surveillance”, which is thought to have single-handedly caused a significant decrease of trust in the federal government. When trust in an institution decreases, people are more likely to turn to a competitor, but with research showing a steady decline of trust in technology companies in light of happenings such as the Cambridge Analytica scandal (Selinger \& Hartzog, 2016), these non-governmental contract-tracking apps are distrusted as well.
\\ \\
This situation of distrust towards both governmental and privately owned contact-tracing apps led to a situation in which the government could chose to use other means in order to track the contagion of the virus, but at the cost of using their power to either commandeer location and interaction information from cellular service providers, or to mandate the use of these contact-tracing apps.
\\ \\
According to (Urbaczewski \& Lee, 2020), the biggest obstacles standing in the way of getting the masses to use contact-tracing apps are ethical concerns related to privacy, security, and anonymity. Their article also points out that many of these technical solutions may have been rolled out without fully understanding the potential future consequences of their actions. Privacy policies, anonymity requirements and data retention policies may not have been at a desirable level before the implementation of the contact-tracing apps, and due to the uniqueness of the situation, legislation surrounding the use of this data is unclear on whether or not other government agencies, not related to the battle against the pandemic, may gain access to this data. 
\\ \\
This is where the differences between separate countries’ handling of this issue begin to arise. Where laws like the EU’s General Data Protection Regulation (GDPR) function to control the use of individual data by companies, it does not specify whether or not these regulations will be of benefit to citizens whose government is collecting data on them. The UK on the other has provided their government with the NHS act of 2006 which allows them to store personally identifying information from contact-tracing apps for a period of 20 years so long as it is used only for the health emergency (Urbaczewski \& Lee, 2020).
\\ \\
Aside from implementing new laws like in the case of the UK, another way to handle this obstacle of public concern is show that these countermeasures are in fact making a difference and thus increasing the people’s level of trust as well as their willingness to infringe upon their own privacy for the sake of the greater good (Ribeiro-Navarrete, et al., 2021) (Hargittau, et al., 2020). The situation of a global pandemic is something that we as a society have not found ourselves in since the outbreak of the Spanish Flu in 1918, and now that we have actually gathered the data to prove that the impact of mobile tracking technology on covid-19 cases per day in countries with mandated use of contact-tracing apps is highly significant to the point where it may decrease the amount of new cases per day by 3.3 on average. While this impact may appear insignificant, it would dramatically flatten the exponential growth of new cases if a policy enforcing the use of these applications were introduced in the early stages of a new wave or outbreak (Urbaczewski \& Lee, 2020).
\\ \\
Yet another argument in favor of this form of mass-tracking is that in the case that the government truly wanted to track a person of interest, it would not be necessary to for that person to consent to downloading a tracking app in order for it to be possible.
As of 2019 over 2/3 of global citizens, and nearly 100 $\%$ of the more highly developed world had access to mobile phones, and in particular smartphones. While in use, mobile phones will constantly register their presence by “pinging” local cellular towers to establish a connection to a network. Through a process called triangulation, it is possible to use take the signals of multiple towers and calculate a mobile phone’s more exact location. What this means for us is that by carrying your phone on your person, the possibility of tracking your location is fully possible (Urbaczewski \& Lee, 2020). This is enforced by the previously discussed concept of companies having common authority over their business records and therefore being able to consent to a government search or acquisition of this data regardless of the users’ wishes. 
While this may be an inefficient way of tracking, only suitable in cases where it is of the utmost importance to gain insight into a person’s whereabouts, it is important to note that through methods like these, there will always be some way to track a mobile phone, and by extension, a person.
\\
\subsection{How powerful is really more data for large corporations?}
The Central Limit Theorem (CLT) is "the most important theorem of probability" (Devore, 2018, p. 298). It origins from the observation that "averaging produces a distribution more bell-shaped than the one being sampled" (Devore, 2018, p. 298). In fact, the mean $\mu$ will be approximately normal for large $n$. As a rule of thumb, "if n > 30, the Central Limit Theorem can be used" (Devore, 2018, p. 302). So if we wish to calculate a probability for the mean in order to create confidence intervals or carry out hypothesis testing, "we need only 'pretend' that the sample mean is normal, standardize it, and use the normal table" (Devore, 2018, p. 299). And this by only requiring a manageable sample size. 
\\ \\
Another branch under data science is machine learning, a science that is about building models algorithmically based on the sample. In statistics, known models are tested, and their performance in relation to the sample is studied. This is not necessary in machine learning where the power of computers are leveraged. It is interesting to see how impactful this approach has become. A quick news search on AI will paint a story of AI supremacy. "AI 'outperforms' doctors diagnosing breast cancer"; "AI system outperforms humans in designing floorplans for microchips"; "AI Outperforms Humans in Question Answering"... It has been able to vastly expand on data science from what statistics has offered. Different from statistics, machine learning arguably thrives more on large sample size. In a publication titled "Do We Need More Training Data?" the authors write (Zhu et. al, 2016) "Based on our analysis, we conjecture that the greatest gains in detection performance will continue to derive from improved representations and learning algorithms that can make efficient use of large datasets," (p. 76) and highlight importance on correctly tuned models over sample size. The research was done on object recognition systems. For earlier conducted research in the field of Natural Language Processing, different evidence was found. "At least for the problem of confusable disambiguation, none of the learners tested is close to asymptoting in performance at the training corpus size commonly employed by the field" (Banko, Michele and Brill, 2001). The term \textit{confusion set disambiguation} is explained to be "the problem of choosing the correct use of a word, given a set of words with which it is commonly confused." The total corpus size in this experiment reached one billion, which is a large data-set indeed and might explain their findings. 
\\
\subsection{Is surveillance moral according to Immanuel Kant's moral philosophy?}
Immanuel Kant believed that we had freedom and a free will to make moral choices because of the distinction he makes "between \textit{phenomena}, which is what we know through experience, and \textit{noumena}, which we can consistently think but not know through experience," and which "allow us to resolve the 'antinomy' about free will by interpreting the 'thesis' that free will is possible as about noumena and the 'antithesis' that every event has a cause as about phenomena" (Johnson, Robert and Cureton, 2022, chs. 1). Trying to make sense of morality and what it is, Kant insists that morality can only be identified \textit{a priori} and not \textit{a posteriori}. In article "Kant’s Moral Philosophy," the authors write "given that the analysis of concepts is an a priori matter, to the degree that ethics consists of such an analysis, ethics is a priori as a well," and also that "ethics is the nature and content of the principles that necessarily determine a rational will" (Johnson, Robert and Cureton, 2022, chs. 1). The "Categorical Imperative" (CI) is the most famous outline of this idea and "Kant’s first formulation of the CI states that you are to 'act only in accordance with that maxim through which you can at the same time will that it become a universal law'" (Johnson, Robert and Cureton, 2022, chs. 5). This puts boundaries on what is morally permissible, or simply moral. Trying to look at this in the context of our thesis, we can note that if surveillance encourages conformity, the positive behavior would not be rooted in a steadfast will. "What makes a good person good is his possession of a will that is in a certain way 'determined' by, or makes its decisions on the basis of, whatever basic moral principles there may be" (Johnson, Robert and Cureton, 2022, chs. 2). Surveillance would minimize the will. It's the freedom to make choices that gives rise to morality. "A \textit{holy} or \textit{divine} will, if it exists, though good, would not be good because it is motivated by thoughts of duty because such a will does not have natural inclinations and so necessarily fulfills moral requirements without feeling constrained to do so" (Johnson, Robert and Cureton, 2022, chs. 2). Conformity would additionally be an act in self-interest, which is amoral because it can't be made universal. Morality as Kant views it can paint a very inflexible and rigid attitude. The philosopher Benjamin Constant made a critique in the form of a thought experiment --- If lying is amoral, would it still be wrong to lie to a murderer who asks for the whereabouts of his intended victim? 
\\ \\

\section{Conclusion}
By looking at what we have already discussed in this article, it is clear to see that we have entered a new era regarding the availability, generation, utilization and commerce of user data. While the original purpose of this new paradigm was to maintain and ensure the safety of nations and their inhabitants, it has brought with it a culture in which user data is treated as a valuable resource and is therefore in high demand, which again motivates companies to collect more data on their users as well as urging them to volunteer more information about themselves.
\\ \\
It can sometimes be difficult to tell where user privacy begins and ends in a situation like the one we currently find ourselves in. While the information that users chose to display on their online profiles is entered knowingly, the amount and variety of information a digital service or application can discern about their user is only limited by what they want to include in their terms of service agreement. Whatever information this application chooses to collect will then be their own through the concept of common authority, it can then be accessed by the U.S. government’s Supreme Court protected right to enter these companies’ databases and business records. In regard to user privacy, the users are likely aware that if they chose to display identifying information like their name, date of birth, former school, political leaning et cetera online, this information is available to the public, and may be used however the finder may please. What the user may however not be aware of is what kind of data is collected without any indication that it is happening, such as geological data of where a picture may have been taken, or what types of content they tend to interact with, which as previously discussed can be used to gain insight into the user’s background, demography, lifestyle, and behavioral patterns. In other words, the user may choose to state their name and age, but the application may discern information about who they are as a person.
\\ \\
Another important thing to consider are the ways in which we as people allow both authorities and private companies to influence our lives by either knowingly or unknowingly sharing our data. It is a fact that applications like Messenger, WhatsApp and TikTok collect information from their users’ messages in order to make personalized advertisements, and that Facebook will show their users advertisements of products they may only have looked at on a website that they perceived as completely unrelated. It is hard to get exact numbers that represent the way this form of targeted advertisement has affected the spending habits of its audience, but judging from the sheer amount of companies that partake in these less traditional forms of advertisement as well as the results we’ve already discussed in this paper, we can likely assume that it has increased the profits of the majority of companies engaging in it.
\\ \\
Aside from targeted advertisements, the users of social platforms such as Facebook also risk making themselves vulnerable to social engineering, manipulation or outright recruitment by potentially malicious parties through the sale and purchase of user data. As previously mentioned, the Cambridge Analytica scandal proved that the data which Facebook collects from its users can potentially be used to change the outcome of political elections. The same data that was used to sway the voters through unique advertisements appealing to the many different categories of people can also be used by malicious groups to help identify users that could be more susceptible to indoctrination or recruitment. 
\\ \\
Furthermore, personal information that we may think is private may turn out to be anything but due to documents like court records, rap sheets or credit statements only being protected by the Supreme Court concept of practical obscurity. This may result in that information being available to the government, landlords, insurance companies and employers through companies that specialize in collecting this information from multiple sources in order to create purchasable aggregate reports of given individuals. Aggregate reports like these could in some cases determine whether a person is considered qualified for a loan, job, insurance, or professional license and thus have great effect on the person’s life.
\\ \\
Another notable point of interest is that while a person may refuse to interact with any form of website or application that collects data about their geographic location, such as contact-tracing apps, through simply carrying their mobile phone with them, government agencies have the possibility of tracking their location and movements. While this may be more of an assurance to the public safety than a cause of worry it is important to note that this solidifies the fact that there are some forms of tracking no amount of reading the terms of service can prevent from affecting you.
\\ \\
To conclude this summary of the way we allow ourselves to be affected by the sharing of our data it seems appropriate to mention that data is something of value, and as is always the case of something of value, there will be parties with motivation to obtain it, legally or illegally, often with the intention of gaining something from the individuals corresponding to the data. There is no way of telling in which ways the data we leave behind every day could be used in the future, but it could be beneficial to be wary of what information we purposefully present to the world as there is power in information, and the chance of it ending up in the wrong hands is ever present in this era of mass data collection, tracking and surveillance.
\\ \\

% Kort oppsummering av hva du har kommet frem til. Oppsummer hvordan oppgaven svarte på motivation/spørsmål fra introduksjonen av oppgaven.


%\bibliography{refs}
%\bibliographystyle{plain}

\section{References}
Devore, J. L., \& Berk, K. L. (2018) Modern Mathematical Statistics with Application. New York Dordrecht Heidelberg London: Springer.
\\ \\
Zhu, X., Vondrick, C., Fowlkes, C.C. et al. Do We Need More Training Data?. Int J Comput Vis 119, 76–92 (2016). https://doi.org/10.1007/s11263-015-0812-2. 
\\ \\
Banko, Michele and Eric Brill. “Scaling to Very Very Large Corpora for Natural Language Disambiguation.” ACL (2001).
\\ \\
Johnson, Robert and Adam Cureton, "Kant’s Moral Philosophy", The Stanford Encyclopedia of Philosophy (Fall 2022 Edition), Edward N. Zalta \& Uri Nodelman (eds.), URL = https://plato.stanford.edu/archives/fall2022/entries/kant-moral/. 
\\ \\
Alava, S., Frau-Meigs, D. & Hassan, G., 2017. Youth and Violent Extremism on Social Media: Mapping the Research. s.l.:United NAtions Educational, Scientific and Cultural Organization.
\\ \\
Arya, V., Sethi, D. & Paul, J., 2019. Does digital footprint act as a digital asset? – Enhancing brand experience through remarketing. s.l.:International Journal of Information Management.
\\ \\
Buchanan, T., Paine, C., Joinson, A. N. & Reips, U., 2006. Development of measures of online privacy concern and protection for use on the Internet. s.l.:Wiley Online Library.
Cate, F. H., 2008. Government data mining: The Need for a Legal Framework. s.l.:Harvard Civil Rights-Civil Liberties Law Review.
\\ \\
Georgieva, I., Beaunoyer, E. & Guitton, M. J., 2021. Ensuring social acceptability of technological tracking in the COVID-19 context. s.l.:Elsevier.
\\ \\
Hargittau, E., Redmiles, E. M., Vitak, J. & Zimmer, M., 2020. Americans' willingness to adopt a covid-19 tracking app. s.l.:University of Zurich.
\\ \\
Karas, S., 2002. Enhancing the privacy discourse: consumer information gathering as surveillance.. s.l.: J. Tech. L. & Pol'y..
\\ \\
Kerr, O., 2021. Buying Data and the Fourth Amendment. s.l.:Lawfare Institute.
\\ \\
Kaashoek, J. & Santillana, M., 2020. COVID-19 positive cases, evidence on the time evolution of the epidemic or an indicator of local testing capabilities? A case study in the United States.. s.l.:SSRN.
\\ \\
McCausland, P. & Schnecter, A., 2018. Trump-Linked Consultants Harvested Data from Millions on Facebook.. s.l.:NBCUniversal News Group.
\\ \\
Mulvenna, M. D., Anand, S. S. & Büchner, A. G., 2000. Personalization on the Net using Web mining. s.l.:Communications of the ACM.
\\ \\
Munzert, S. et al., 2021. Tracking and promoting the usage of a COVID-19 contact tracing app. s.l.:Nature Human Behaviour.
\\ \\
Palos-Sanchez, P., Saura, J. R. & Martin-Velicia, F., 2019. A study of the effects of programmatic advertising on users' concerns about privacy overtime. s.l.:Elsevier.
\\ \\
Personal Finance, 2019. What is a ChoicPoint Rrport?. [Internet] 
Available at: https://personal-finance.extension.org/what-is-a-choicepoint-report/
[Found 25. October 2022].
\\ \\
Rathi, R., 2019. Effect of Cambridge Analytica’s Facebook ads on the 2016 US Presidential Election. s.l.:Towards Data Science.
\\ \\
Reporters Committee for Freedom of the Press, 2009. Out of sight, out of bounds: Article from the Spring 2009 issue of The News Media & The Law. [Internet] 
Available at: https://www.rcfp.org/journals/the-news-media-and-the-law-spring-2009/out-sight-out-bounds/
[Found 25. October 2022].
\\ \\
Ribeiro-Navarrete, S., Saura, J. R. & Palacios-Marqués, D., 2021. Towards a new era of mass data collection: Assessing pandemic surveillance technologies to preserve user privacy.. s.l.:National Library of Medicine.
\\ \\
Russel, M. A., 2013. Mining the Social Web. 2. red. s.l.:O'Reilly Media.
\\ \\
Selinger, E. & Hartzog, W., 2016. Facebook’s emotional contagion study and the ethical problem of co-opted identity in mediated environments where users lack control. s.l.:Sage Journals.
\\ \\
Sinha, G. A., 2013. NSA surveillance since 9/11 and the human right to privacy. s.l.:s.n.
\\ \\
Sumpter, D. J. T., 2018. David J. T. Outnumbered: from Facebook and Google to Fake News and Filter-Bubbles — the Algorithms That Control Our Lives.. s.l.:Bloomsbury Sigma.
\\ \\
Testa-Avila, E., 2008. Practical Obscurity in the Digital Age: Public Records in the Private Sector. s.l.:UC Berkeley.
\\ \\
Trang, P. T., 2017. Personalized ads on Facebook: An effective marketing tool for online marketers. s.l.:Elsevier.
\\ \\
Urbaczewski, A. & Lee, Y. J., 2020. Information Technology and the pandemic: a preliminary multinational analysis of the impact of mobile tracking technology on the COVID-19 contagion control. s.l.:Taylor & Francis Online.
\\ \\
Vargo, C. J. & Hopp, T., 2020. Fear, Anger, and Political Advertisement Engagement: A Computational Case Study of Russian-Linked Facebook and Instagram Content. s.l.:Journalism & Mass Communication Quarterly.
\\ \\
Zuboff, S., 2015. Big other: surveillance capitalism and the prospects of an information civilization. s.l.:Journal of Information Technology.
\\ \\
Zuboff, S., 2019. The Age of Surveillance Capitalism: The Fight for a Human Future at the New Frontier of Power. s.l.:Profile Books.

\end{document}
